%MIT OpenCourseWare: https://ocw.mit.edu
%RES.18-012 Algebra II Student Notes, Spring 2022
%License: Creative Commons BY-NC-SA 
%For information about citing these materials or our Terms of Use, visit: https://ocw.mit.edu/terms.

% \rhead{\rightmark}
\lhead{Lecture 27: Finite Fields}

\section{Finite Fields}

Previously, we constructed the finite field $\mathbb{F}_q$ for $q = p^n$, and showed that there is a unique such field. This construction had the unusual property that the field \emph{consisted} of exactly the roots of a polynomial (where the polynomial was $x^q - x$), rather than just being \emph{generated} by the roots of a polynomial. 

There is more that we can say about the structure of finite fields. 

\subsection{The Multiplicative Group}

\vspace{0.75em}

\begin{lemma}\label{cyclicmultgroup}
    If $F$ is any field and $G$ is a finite subgroup of $F^\times$, then $G$ is cyclic. 
\end{lemma}

\begin{example}
    If $F = \mathbb{C}$, then finite subgroups of $F^\times$ are the $n$th roots of unity \[\left\{\exp \frac{2\pi i}{n} \right\} = \langle \zeta_n \rangle \cong \mathbb{Z}/n.\] 
\end{example}

\begin{proof}[Proof of Lemma \ref{cyclicmultgroup}]
    By the classification of finite abelian groups, we know $G \cong \prod \mathbb{Z}/p_i^{n_i}\ZZ$ for some integers $n_i$. So it's enough to check that no prime appears twice (meaning that every prime appears in the list of $p_i$ at most once). Then we can use the Chinese Remainder Theorem to show that the product is cyclic, as all the $p_i$ are then coprime. For example, $\mathbb{Z}/4\ZZ \times \mathbb{Z}/3\ZZ = \mathbb{Z}/12\ZZ$ is cyclic, while $\mathbb{Z}/4\ZZ \times \mathbb{Z}/2\ZZ$ is not. 

    But suppose $p$ appears twice. Then $G$ contains a subgroup $\mathbb{Z}/p^a\ZZ \times \mathbb{Z}/p^b\ZZ$. But then $G$ has at least $p^2$ elements of order dividing $p$, since there's $p$ choices for the coordinate in each. This would mean the polynomial $x^p - 1$ has at least $p^2$ roots; but it has degree $p$, so this is impossible. 
\end{proof}

\begin{corollary}
    For any finite field $\mathbb{F}_q$, its multiplicative group $\mathbb{F}_q^\times$ is cyclic, meaning $\mathbb{F}_q^\times \cong \mathbb{Z}/(q - 1)$. 
\end{corollary}

\begin{note}
    Although we know in theory that $\mathbb{F}_q^\times \cong \mathbb{Z}/(q - 1)$, in practice it's hard to compute how this isomorphism works --- it is difficult to find a generator, or to figure out what power to raise the generator to in order to get a given element. Many cryptography and encryption protocols are based on this. 
\end{note}

\begin{corollary}
    We have $\mathbb{F}_q \cong \mathbb{F}_p(\alpha)$, and therefore, there exists an irreducible polynomial of any degree over $\mathbb{F}_p$. 
\end{corollary}

\begin{proof}
    There exists $\alpha \in \mathbb{F}_q$ which generates the multiplicative group; then $\alpha$ must generate $\mathbb{F}_q$ as an extension of $\mathbb{F}_p$, since every element of $\mathbb{F}_q$ is a \emph{power} of $\alpha$. (The converse is false --- it is possible to find $\alpha$ which generates the extension but not the multiplicative group.) 

    Then $\mathbb{F}_q = \mathbb{F}_p(\alpha) \cong \mathbb{F}_p[x]/(Q)$ where $Q$ is the minimal polynomial of $\alpha$. So $Q$ is an irreducible polynomial of degree $n$, where $q = p^n$. In particular, a procedure (in theory) to find an irreducible polynomial of degree $n$ is to write down $\mathbb{F}_{p^n}$, find a multiplicative generator, and take its minimal polynomial.
\end{proof}

\subsection{Application to Number Theory}

Finite fields arise in many areas of math and computer science --- in particular, in number theory. One such example is $R/(p)$, where $R$ is the ring of algebraic integers in a finite extension of $\mathbb{Q}$. 

\begin{example}
    If $p \equiv 3 \pmod{4}$, then $\mathbb{Z}[i]/(p) \cong \mathbb{F}_{p^2}$ --- it's a field since $(p)$ is maximal. 
\end{example}

The example we'll focus on is the extension $\mathbb{Q}(\zeta_\ell)$, where $\ell$ is a prime (it's possible to consider general $\ell$, but the prime case is a bit simpler). We know this extension is $\mathbb{Q}[x]/(x^{\ell - 1} + \cdots + 1)$, and we can check that its ring of algebraic integers is \[R = \mathbb{Z}[x]/(x^{\ell - 1} + \cdots + 1).\] 

\begin{qq}
For an (integer) prime $p$, when is $R/(p)$ a field (or equivalently, when is $(p)$ maximal)?
\end{qq}

It's clear that $R/(p) \cong \mathbb{F}_p[x]/(x^{\ell - 1} + \cdots + 1)$, which means its dimension over $\mathbb{F}_p$ (as a vector space) is $\ell$. So if $R/(p)$ is a field, then it must be $\mathbb{F}_{p^\ell}$. 

We'll assume $p \neq \ell$. 

\begin{proposition}
    If $p \neq \ell$, then $R/(p)$ is a field if and only if $\ord_{\mathbb{F}_\ell^\times} p = \ell - 1$.
\end{proposition}

Here $\ord_{\mathbb{F}_\ell^\times} p$ denotes the multiplicative order of $p$ mod $\ell$; in particular, $\ell - 1$ is the largest possible order, since $\mathbb{F}_\ell^\times$ has $\ell - 1$ elements. 

\begin{proof}
    Let $\mathfrak{m}$ be a maximal ideal of $R$ containing $(p)$. Then we have $R/\mathfrak{m} \cong \mathbb{F}_{p^a}$ for some $a$. Let the image of $\zeta_\ell$ in $R/\mathfrak{m}$ be $\overline{\zeta}_\ell$. Then we know $\overline{\zeta}_\ell^\ell = 1$ and therefore $\overline{\zeta_\ell}$ has multiplicative order $\ell$; so since the multiplicative group of  $\mathbb{F}_{p^a}$ has size $p^a - 1$, we get that $\ell \mid p^a - 1$. 

    Now if the order of $p$ in $\mathbb{F}_\ell^\times$ is $\ell - 1$, then we must have $a \geq \ell - 1$. But it cannot be larger than $\ell - 1$. So then $a = \ell - 1$ and $R/\mathfrak{m} \cong R/(p)$, which means $R/(p)$ is a field. The converse can be proved similarly. 
\end{proof}

\begin{example}
    Suppose $p = 3$ and $\ell = 5$. Then $\ord_5 3 = 4$, so $R/(3)$ is a field. 
\end{example}

\subsection{Multiple Roots}

In our construction of finite fields, one step had to do with multiple roots and derivatives. In particular, we used the fact that a multiple root of $P$ is also a root of $P'$ in order to show that the Artin--Schreier polynomial doesn't have multiple roots. 

\begin{qq}
    Let $P \in F[x]$ be an irreducible polynomial. Can $P$ have multiple roots in its splitting field (or equivalently, in any extension)?
\end{qq}

If $\alpha$ is such a root, then $\alpha$ is also a root of $P'$, and therefore a root of $\gcd(P, P')$ as well (where $\gcd(P, P')$ is the polynomial $Q$ which generates $(P, P')$ as an ideal). 

But $P$ is irreducible, and $\deg P' < \deg P$. So if $P' \neq 0$, then this means $\gcd(P, P') = 1$, and no such $\alpha$ can exist. However, it's possible that $P' = 0$. So the question reduces to the following:

\begin{qq}
    When can we have a nonconstant polynomial with $P' = 0$?
\end{qq}

We have $(x^n)' = nx^{n - 1}$. If $n \geq 1$, then if the field has characteristic $0$, this is always nonzero. Meanwhile, if the field has characteristic $p$, then this is zero if and only if $p \mid n$. So if $P' = 0$, then we must have \[P(x) = Q(x^p) = a_nx^{pn} + a_{n - 1}x^{p(n - 1)} + \cdots + a_0,\] where $p = \operatorname{char}(F)$. So we want to see when such a polynomial is irreducible. 

If $F = \mathbb{F}_q$ is finite, then we know $a^q = a$ for all $a \in F$. This means we can extract $p$th roots of the coefficients, since $(a^{p^{n - 1}})^p = a$ --- so we can write $a_i = b_i^p$ for some $b_i \in F$. Then we have \[P(x) = b_n^px^{pn} + b_{n - 1}^px^{p(n - 1)} + \cdots + b_0^p.\] But this allows us to extract a $p$th root of the \emph{polynomial}: we then have \[P = (b_nx^n + b_{n - 1}x^{n - 1} + \cdots + b_0)^p.\] 

On the other hand, there exist examples of such irreducible $P$ in infinite fields. For instance, take $F = \mathbb{F}_q(t)$ to be the field of rational functions in $t$ (or equivalently, the fraction field of $\mathbb{F}_q[t]$), and $P(x) = x^p - t$. This is irreducible, but its derivative is identically $0$. 

We won't study examples like this, but it's good to know they exist --- in every situation we care about, irreducible polynomials can't have multiple roots.

\begin{definition}
    An extension $E/F$ is \textbf{separable} if the minimal polynomial (over $F$) of every algebraic element $\alpha \in E$ has no multiple roots. 
\end{definition}

So if $F$ has characteristic $0$ or is finite, then every extension is separable. We'll only look at these instances, so we will generally assume all our extensions are separable. 

\subsection{Geometry of Function Fields}

Another important example of a field is $\mathbb{C}(t)$; we can think of the extensions of $\mathbb{C}(t)$ via geometry. Let $F = \mathbb{C}(t)$, and suppose $E = F[x]/(P)$ is a finite extension of $F$, where $P$ is an irreducible polynomial. As with integers, we can scale so that $P$ is a primitive polynomial in $\CC[t][x]$. 

We can then think of $P$ as a polynomial in two variables, meaning $P \in \mathbb{C}[t, x]$. So another way to think of these extensions is that $F = \operatorname{Frac}(\mathbb{C}[t])$, and $E = \operatorname{Frac}(R)$ where $R = \mathbb{C}[t, x]/(P)$.

As discussed earlier, these rings are worked with in algebraic geometry --- to connect them to geometry, we consider the maximal spectrum \[X = \operatorname{MSpec}(R) = \{(a, b) \in \mathbb{C}^2 \mid P(a, b) = 0\}\] (which describes all maximal ideals of $R$). Also define $Y = \operatorname{MSpec}(\mathbb{C}[t]) = \mathbb{C}$. Then we have a map $X \to Y$ sending $(a, b) \mapsto a$. 

\begin{question}
    If $P$ is irreducible, is $R$ a field extension?
\end{question}

\begin{ans}
No, $R$ is not a field. Polynomials in two variables aren't a PID (so even if $P$ is irreducible, $(P)$ is generally not maximal) --- if they were, algebraic geometry would be trivial. 
\end{ans}

\begin{example}
    Let $P(t, x) = x^n - t$. 
\end{example}

Then $R \cong \mathbb{C}[x]$, where this map sends a complex number to its $n$th power (since $t = x^n$). Each point in $\mathbb{C}$ has $n$ complex $n$th roots (except $0$), giving a \emph{ramified covering} (with a ramification point at $0$). One way to represent this geometrically is to draw the $t$-plane and the $x$-plane. In the $t$-plane, we make a cut along the $x$-axis, turning it into two half-planes glued together. 

\begin{center}
    \begin{asy}
        unitsize(2cm);
        dot((0, 0));
        draw((1, 0)--(-1, 0), Arrows);
        draw((0, 1)--(0, -1), Arrows);
        fill((1, 0)--(1, 1)--(-1, 1)--(-1, 0)--cycle, heavycyan+opacity(0.1));
        fill((1, 0)--(1, -1)--(-1, -1)--(-1, 0)--cycle, heavyred+opacity(0.1));
        label("$t$", (0, 1), dir(90));
    \end{asy}
\end{center}

For a point on the $x$-plane, raising it to the $n$th power multiplies the angle by $n$. So we cut the $x$-plane into $2n$ pieces (colored by which half-plane their points are mapped to):

\begin{center}
    \begin{asy}
        unitsize(2cm);
        dot((0, 0));
        draw((1, 0)--(-1, 0), Arrows);
        draw((0, 1)--(0, -1), Arrows);
        draw(dir(60)--dir(240), gray);
        draw(dir(120)--dir(300), gray);
        fill((0, 0)--(1, 0)--dir(60)--cycle, heavycyan+opacity(0.1));
        fill((0, 0)--dir(60)--dir(120)--cycle, heavyred+opacity(0.1));
        fill((0, 0)--dir(120)--(-1, 0)--cycle, heavycyan+opacity(0.1));
        fill((0, 0)--(-1, 0)--dir(240)--cycle, heavyred+opacity(0.1));
        fill((0, 0)--dir(240)--dir(300)--cycle, heavycyan+opacity(0.1));
        fill((0, 0)--dir(300)--(1, 0)--cycle,  heavyred+opacity(0.1));
        label("$x$", (0, 1), dir(90));
    \end{asy}
\end{center}

This describes the geometry of the map raising $x$ to the $n$th power. 

\begin{example}
    Let $P(x, t) = x^2 - t(t + 1)(t - \lambda)$. (Any $P$ consisting of $x^2$ minus a cubic polynomial can be written in this form, by a change of variables.) For simplicity, assume $\lambda \in \mathbb{R}$.
\end{example}

We can again draw the $\mathbb{C}$-plane. We again have a ramified double covering, with three ramification points --- $t = 0$, $-1$, and $\lambda$ (for every other point, there are two square roots). So we can again make a cut and create two half-planes.  

\begin{center}
    \begin{asy}
        unitsize(2cm);
        dot((0, 0));
        dot("$-1$", (-0.4, 0), dir(270));
        dot("$\lambda$", (0.6, 0), dir(270));
        draw((1, 0)--(-1, 0), Arrows);
        draw((0, 1)--(0, -1), Arrows);
        fill((1, 0)--(1, 1)--(-1, 1)--(-1, 0)--cycle, heavycyan+opacity(0.1));
        fill((1, 0)--(1, -1)--(-1, -1)--(-1, 0)--cycle, heavyred+opacity(0.1));
        label("$t$", (0, 1), dir(90));
    \end{asy}
\end{center}

For each half-plane, its pre-image breaks into two pieces (corresponding to the two branches of the square root --- we can start with the positive or negative one). So the pre-image consists of two blue rectangles and two red rectangles:

\begin{center}
    \begin{asy}
    unitsize(1.4cm);
    fill((0, 0)--(1, 0)--(1, 1)--(0, 1)--cycle, heavycyan+opacity(0.1));
    fill((1.5, 0)--(2.5, 0)--(2.5, 1)--(1.5, 1)--cycle, heavycyan+opacity(0.1));
    fill((0, -0.5)--(1, -0.5)--(1, -1.5)--(0, -1.5)--cycle, heavyred+opacity(0.1));
    fill((1.5, -0.5)--(2.5, -0.5)--(2.5, -1.5)--(1.5, -1.5)--cycle, heavyred+opacity(0.1));
    \end{asy}
\end{center}

We then need to glue these rectangles together, by thinking about the values of these functions. When glued together, they look like a bagel (where we cut the bagel horizontally and through its middle). 

\begin{note}
These situations require more background to describe rigorously, and for that reason they are usually not presented in algebra classes; but they are important examples of field extensions, and mathematicians often have these examples in mind even when constructing algebraic arguments about number fields. 
\end{note}